% STATUS: conditional framework; outstanding model-specific work is explicit.
\section{P18: Stability radius and Hopf shift (unfinished application)}
\label{sec:P18}
\subsection{Moving Jacobian}\label{slot:P18-1}
For P15's specified family set
$E(\eps)=D_xF(x^*(\eps),\eps)-J_0$. Bounds must include both direct
coefficient variation and equilibrium motion; these estimates are absent
from the supplied outline.
\subsection{Sufficient stability radius}\label{slot:P18-2}
If $J_0$ is normal and Hurwitz, put $\gamma=-\sa(J_0)>0$.
The normal-matrix perturbation inclusion
\cite[\S1 and Theorem~$4_\eps$]{EmbreeTrefethen2001} guarantees
stability for $\|E(\eps)\|_2<\gamma$, even when the perturbed matrix
is nonnormal. Turning this into a physical parameter radius requires
an explicit bound on $E$; it is not a sharp structured-radius assertion.
\subsection{Simple eigenvalue and threshold shift}\label{slot:P18-3}
Separately, at a simple transverse Hopf point of an analytic family,
\cite[Theorem~1]{GreenbaumLiOverton2020} and the implicit-function
theorem give
\[
 g_H'(0)=-\frac{\Ree[p^*J_\eps v]}{\Ree[p^*J_gv]},\qquad
 p^*v=1,\quad\Ree[p^*J_gv]\ne0.
\]
Here $J_\eps,J_g$ differentiate the Jacobian evaluated at the moving
equilibrium. These are not derivatives with the equilibrium frozen.
\subsection{Remaining Hopf conditions}\label{slot:P18-4}
Verify the family's regularity, remaining spectral gap, crossing and
first-Lyapunov-coefficient nondegeneracy. A spectral threshold shift alone
does not establish a nonlinear Hopf bifurcation.
